\noindent\textbf{(a)} $x^2+y^2=C$, $C>0$; como gráficas, $y=\pm\sqrt{C-x^2}$ en $|x|<\sqrt C$.

\noindent\textbf{(b)} $y=2x$, en el intervalo máximo $(0,\infty)$ del PVI.

\noindent\textbf{(c)} $(x+1)y^2+x^3/3=C$. Las gráficas ordinarias son $y=\pm\sqrt{(C-x^3/3)/(x+1)}$ en intervalos con $x\ne-1$ y radicando positivo.

\noindent\textbf{(d)} $e^{xy}+\sin x=C$. Para $x\ne0$, $y=\log(C-\sin x)/x$ con $C-\sin x>0$. Sólo $C=1$ admite prolongación por $x=0$, con $y(0)=-1$ y $y^{\prime}(0)=-1/2$.

\noindent\textbf{(e)} $x^2+xy-2y^2=C$. Para $C\ne0$, $y=(x\pm\sqrt{9x^2-8C})/4$ donde el radicando es positivo. El nivel $C=0$ da las rectas $y=x$ y $y=-x/2$.

\noindent\textbf{(f)} $y=\log(C-e^x)$, $C>0$, $x<\log C$.

\noindent\textbf{(g)} $x^3+3xy+y^3=C$, interpretado por sus arcos gráficos $C^1$. Donde $x+y^2\ne0$ son ramas regulares. El nivel $C=0$ posee una rama suave por el origen de forma $y=-x^2/3+O(x^5)$; el punto $(-1,-1)$ del nivel $C=1$ es aislado localmente y no una trayectoria.
